\documentclass[11pt]{article}
\usepackage[margin=1in]{geometry}
\usepackage{amsmath,amssymb}
\usepackage{booktabs}
\usepackage[hidelinks]{hyperref}
\usepackage{natbib}
\usepackage{graphicx}

\title{Computing Through Tools Decorrelates Same-Model LLM Ensembles}
\author{Reda Louriki\\\texttt{lourikireda06@gmail.com}}
\date{}

\begin{document}
\maketitle

\begin{abstract}
Agents built on the same large language model (LLM) share weights, so repeated sampling or
persona prompts tend to produce correlated errors, which limits the benefit of ensembling them.
We study how to induce useful error diversity among same-model agents and when that diversity
compensates for the individual accuracy it costs. Applying the ambiguity decomposition to the
vote distribution of an agent ensemble gives an exchange-rate diagnostic that is exact for the
Brier score of the vote distribution; whether it also predicts majority-vote outcomes is tested
empirically. On recent competition mathematics problems from the MathArena stream, we compare
five ways of diversifying five agents of the same model: repeated sampling, persona prompts,
fixed method templates, planner-generated plans, and computational substrates, in which agents
solve in plain text, executed Python, SymPy, brute-force search, or solve-then-check code.
Pooled across <<NMODELS>> models and <<SA_N>> paired items, substrate ensembles improve
majority-vote accuracy over self-consistency by <<SA_DIFF>> points (95\% CI <<SA_CI>>,
Holm-adjusted $p=<<SA_PHOLM>>$, random effects <<SA_RE>>) without generating more tokens. A
homogeneous control on one model shows that five copies of the best single tool perform as well
as the mixed ensemble (<<HOMOG_ABS>>), and that five identical Python agents are much less
error-correlated than five text agents. The gain therefore comes mainly from computing through
tools rather than from mixing them, although a mixed ensemble reaches the accuracy of the best
tool without requiring that tool to be identified in advance. The effect depends on capability.
Tools raise per-agent accuracy on weak models and mainly decorrelate errors on stronger ones, and
the gain decreases as base accuracy rises (<<CURVE>> points from the weakest to the strongest
model). Planner-generated plans reduce individual accuracy without a significant ensemble
benefit, and the diagnostic anticipates this out of sample.
\end{abstract}

\section{Introduction}
The error of an averaged predictor equals the average member error minus a diversity term
\citep{krogh1995neural,wood2023unified}. For LLMs the simplest ensemble is self-consistency,
which samples one model several times and takes a majority vote \citep{wang2023selfconsistency}.
Its members share weights and training data, and their errors are strongly correlated. A recent
capability-controlled audit reports that common diversity metrics for LLM ensembles mostly track
member weakness and that majority vote rarely beats the strongest member
\citep{audit2026diversity}, and mixing different models can underperform aggregating a single
strong one \citep{li2025rethinking}.

The audit examines whether measured diversity predicts ensemble gains. We examine how useful
diversity can be induced among agents of the same model, and when it pays for the individual
accuracy it costs. Under a fixed budget of five agents per item, we compare stochastic sampling,
surface-level prompt variation, procedural variation through method templates, model-generated
variation through a planner, and substrate variation, where the same model computes through
different tools. Combining chain-of-thought and program-based reasoning has been studied before.
XoT switches between them on verification failure \citep{liu2023plan}, and collaborative
verification combines them with learned verifiers \citep{liang2024improving}. Here we treat
substrates as a source of ensemble diversity, quantify the trade-off between the diversity they
add and the accuracy they cost, separate that diversity from tool augmentation with
homogeneous-substrate controls, and compare substrates with the other mechanisms at equal budget.

\paragraph{Contributions.}
\begin{itemize}
\item An application of ensemble ambiguity theory to LLM agent ensembles on the simplex of
  answer-equivalence classes, giving an exchange-rate diagnostic that is exact for the Brier
  score of the vote distribution and is tested empirically on majority vote
  (Section~\ref{sec:framework}).
\item A controlled comparison of five diversity mechanisms at a fixed agent budget on recent
  competition problems, with Holm-corrected, model-stratified and random-effects analyses.
\item Evidence across four models that tool-substrate ensembles improve over self-consistency
  without generating more tokens, and homogeneous controls indicating that the gain comes mainly
  from computing through tools rather than from heterogeneity itself.
\item A capability-dependent mechanism (augmentation for weak models, decorrelation for stronger
  ones) and a negative result for planner-generated plans.
\end{itemize}

\section{Related work}
\paragraph{Sampling-based ensembles.} Self-consistency \citep{wang2023selfconsistency} and
universal self-consistency \citep{chen2023universal} aggregate samples of one model.
Mixture-of-Agents layers aggregators over several models \citep{wang2024moa}, while same-model
aggregation can perform better \citep{li2025rethinking}.
\paragraph{Prompt diversity.} DIV-SE elicits and ensembles diverse reasoning approaches
\citep{naik2023diversity}, and Dipper ensembles diverse prompts \citep{dipper2024}. We test both
fixed templates and planner-generated approaches while controlling for individual accuracy.
\paragraph{Programs and tools.} Program-aided reasoning offloads computation to an interpreter
\citep{gao2023pal,chen2022program}. XoT integrates chain-, program- and equation-of-thought with
verification-based switching \citep{liu2023plan}, and collaborative verification combines CoT and
PoT with verifiers \citep{liang2024improving}. We study such substrates as ensemble members and
measure their contribution to diversity separately from their contribution to accuracy.
\paragraph{Ensemble diversity theory.} Our accounting instantiates known decompositions
\citep{krogh1995neural,wood2023unified}. Generalized ambiguity decompositions exist for
classification \citep{jiang2017generalized}, and 0/1 loss has no equally clean additive form
\citep{wood2023unified}, which is why our diagnostic is stated for the Brier score of the vote
distribution. Pairwise diversity measures follow \citet{kuncheva2003measures}.

\section{Framework}\label{sec:framework}
Let $K$ agents answer an item with gold answer $y$, with answers mapped to equivalence classes
as described in Section~\ref{sec:setup}. Let $p_i\in\{0,1\}^M$ be the one-hot class of agent $i$,
$\bar p=\frac1K\sum_i p_i$ the vote distribution, and $e_y$ the gold class. Then
\begin{equation}
\underbrace{\tfrac1K\textstyle\sum_i\|p_i-e_y\|^2}_{\text{mean individual error}}
=\underbrace{\|\bar p-e_y\|^2}_{\text{Brier error of the vote distribution}}
+\underbrace{\tfrac1K\textstyle\sum_i\|p_i-\bar p\|^2}_{\text{ambiguity}} .
\end{equation}
Since $\|p_i-e_y\|^2=2\cdot\mathbb 1[\text{agent } i\text{ wrong}]$, a diversity mechanism that
lowers mean accuracy by $\delta$ and raises ambiguity by $\alpha$ improves the Brier error of the
vote distribution if and only if $\alpha>2\delta$. We call $r=\alpha/(2\delta)$ the exchange
rate, and say that a mechanism pays when $r>1$ or $\delta\le 0$.

Majority vote is the argmax of $\bar p$, and 0/1 loss admits no equally clean additive
decomposition \citep{wood2023unified}. Whether $r$ predicts majority-vote gains is therefore an
empirical question, which we test out of sample in Section~\ref{sec:ablation}. We also report
the mean pairwise error correlation $\bar\phi$ and the effective number of independent agents
$K/(1+(K-1)\bar\phi)$.

\section{Experimental setup}\label{sec:setup}
\paragraph{Data.} We use final-answer problems from AIME 2025 and 2026, HMMT (February 2025,
November 2025, February 2026), SMT 2025, CMIMC 2025 and BRUMO 2025, taken from the MathArena
evaluation stream \citep{balunovic2025matharena}. APEX is excluded because it is close to
unsolvable by design. All analyses use one fixed subset of 90 items. Contamination risk varies
by model and competition date, and gemma-4 was released after part of the 2025 competitions.
Older benchmarks were rejected after a probe: MATH levels 4 and 5 are at ceiling for these
models (94 to 100\%), and AIME 1983 to 2024 shows signs of memorisation (90\% accuracy with about
1.7k output tokens).
\paragraph{Answer equivalence.} Final answers are read from the last \verb|\boxed{}| for text
agents and from the last \verb|ANSWER:| line printed by the program for code agents. Two answers
belong to the same class if their normalised strings coincide (after removing \verb|\left| and
\verb|\right|, \verb|\text{}|, degree signs, spaces and trailing periods, and mapping \verb|*I|
to \verb|i|), if \texttt{math-verify} finds them symbolically equivalent in either direction
after LaTeX and expression extraction, or, for plain program output such as \verb|sqrt(12)|,
\verb|pi| or tuples, if the answer converted to LaTeX with SymPy (unevaluated \verb|sympify|)
is found equivalent. Missing or unparseable answers abstain and do not vote. Classes are built
jointly over all answers to an item, including the gold answer, so every system and aggregator
votes over the same labels. All code is released.
\paragraph{Models.} We use gemma-4-31B-it (non-reasoning, NVIDIA NIM), gemini-3.5-flash-lite and
gemini-3.1-flash-lite (Google AI Studio), and gpt-oss-20b (reasoning, served by NVIDIA NIM and by
Groq on disjoint items, so the A and S agents of any item share a provider). Temperature is 0.7
and the output budget is 16k tokens.
\paragraph{Arms.} Each arm has five agents of the same model. \textbf{A} uses an identical
neutral prompt. \textbf{C} uses five fixed method templates. In \textbf{G}, a planner writes five
problem-specific plans and each solver follows one of them. \textbf{S} uses five substrates:
text chain-of-thought, executed Python, SymPy formalisation, brute-force or numeric search, and
solve-then-check. Code runs in a network-isolated sandbox with one repair round. \textbf{X} adds
five more neutral samples, so that A and X together form ten-sample self-consistency.
\textbf{HB} and \textbf{HP} are homogeneous controls with five brute-force or five Python agents.
\paragraph{Aggregation and statistics.} We use majority vote over equivalence classes, with
ties counted at their expected value, cross-fitted weighted votes, and a verifier in which the
base model must pick one of the candidates. Within a model we use a paired item-level bootstrap.
Pooled contrasts use a model-stratified sign-flip permutation test, a stratified bootstrap and
DerSimonian and Laird random effects, with Holm correction over the planned family
$\{S{-}A, C{-}A, G{-}A\}$. Choices made after seeing a model's data, such as removing SymPy
(``S4''), are validated only on the other models. Each item has a single sampled ensemble per
arm, so the item-level bootstrap already includes sampling noise; Section~\ref{sec:stability}
reports a run-to-run replicate.

\section{Results}\label{sec:results}
\subsection{Pilot on GSM8K}
On GSM8K with llama-3.2-11B (100 items), a blind judge recovered the assigned method in 79\% of
template-arm solutions, against 17\% by chance. Templates halved the error correlation
($\bar\phi$ from 0.49 to 0.26) but cost 9.6 points of individual accuracy, an exchange rate of
0.66, and majority-vote accuracy did not change. Persona prompts lowered accuracy without a
significant reduction in error correlation.

\subsection{Main comparison on gemma-4-31B}
\begin{table}[h]\centering\small
\begin{tabular}{lcccccc}\toprule
Arm & Indiv.\ acc & Majority & Verifier & Oracle & $\bar\phi$ & Eff.\ $N$\\\midrule
A self-consistency & 0.800 & 0.867 & 0.889 & 0.922 & 0.575 & 1.52\\
C fixed templates  & 0.793 & 0.867 & 0.900 & 0.956 & 0.443 & 1.80\\
G generated plans  & 0.716 & 0.833 & 0.822 & 0.900 & 0.459 & 1.76\\
S substrates       & \textbf{0.816} & \textbf{0.900} & \textbf{0.922} & \textbf{0.967} & \textbf{0.392} & \textbf{1.95}\\\bottomrule
\end{tabular}
\caption{gemma-4-31B, five agents per arm, the same 90 items in every arm.}
\end{table}
Relative to A, the exchange rate is 3.1 for C and 0.56 for G, and S loses no individual
accuracy.

\subsection{Homogeneous-substrate controls}
<<HOMOG>>

\subsection{Pooled analysis across models}
<<POOLED_TABLE>>
<<TOKENS_TABLE>>

\subsection{Augmentation and decorrelation}
<<MECH_TEXT>>

\subsection{Compute-matched comparison and total cost}
\begin{table}[h]\centering\small
\begin{tabular}{lcc}\toprule
System (gemma-4-31B) & Accuracy & Output tokens per problem\\\midrule
A$\times$5, majority vote & 0.871 & 8.1k\\
A$\times$10 with verifier (compute-matched) & 0.894 & 16.9k\\
S$\times$5 with verifier & \textbf{0.918} & 10.4k\\\bottomrule
\end{tabular}
\caption{Verifier systems on the 85 items for which all verifiers are available.}
\end{table}
<<COMPUTE>>

\subsection{Run-to-run stability}\label{sec:stability}
<<REPL>>

\subsection{Ablation and the exchange-rate diagnostic}\label{sec:ablation}
Leaving one substrate out on gemma shows the largest contributions from brute-force search (3.3
points lost when removed) and executed Python (2.2 points). Removing SymPy looked beneficial in
sample (+1.1 points) but is neutral on held-out models (<<S4OOS>>), so we keep the five-substrate
arm. When the exchange rate is estimated on one half of the items and checked on the other, it
predicts the failure of generated plans in 99\% of random splits for the Brier score, for which
it is exact, and in 84\% of splits for majority vote. For small gains (below about 3 points with
45 items) its predictions are at chance level. The diagnostic is therefore informative for
majority vote, although less reliable than for the quantity from which it is derived.

\subsection{Comparison with the MathArena leaderboard}
MathArena reports single-attempt accuracy of single models, whereas our systems are multi-agent
ensembles at inference time, so the two are not directly comparable and we report output tokens
alongside. On the 85 problems shared by all systems, gemma-4-31B with substrate agents and a
verifier reaches 0.918 with 10.4k output tokens per problem. On the problems they share with
us, o3 (high) scores 0.934 with 9.6k tokens and GPT-OSS-120B (high) scores 0.922. Frontier models
remain 5 to 7 points higher.

\section{Limitations}
Sample sizes are modest, with at most 90 items per model, and gpt-oss-20b is only partially
covered, so its results are supporting evidence. The study covers a single domain, competition
mathematics, and uses free-tier APIs only. Contamination risk varies by model and competition
date. The homogeneous controls were run on one model, so we cannot tell whether heterogeneity
helps on weaker models, where per-tool accuracies differ more. Cost accounting covers LLM input
and output tokens and sandbox time, and includes hidden reasoning tokens only to the extent that
each provider reports them.

\section{Conclusion}
Letting agents of the same model compute through tools is a practical way to induce useful error
diversity. Tools help weak models by raising per-agent accuracy and help stronger models by
decorrelating their errors, without increasing the number of generated tokens. On the model where
we could test it, mixing different tools did not improve on the best single tool replicated five
times; its advantage is that it reaches that level without knowing which tool is best, and the
best tool differs across models. Planner-generated diversity did not help in our experiments. An
exchange-rate diagnostic derived from the ambiguity decomposition identifies when induced
diversity fails to compensate for its accuracy cost.

\bibliographystyle{plainnat}
\begin{thebibliography}{99}
\bibitem[Krogh and Vedelsby(1994)]{krogh1995neural} A.~Krogh and J.~Vedelsby. Neural network ensembles, cross validation, and active learning. In \emph{Advances in Neural Information Processing Systems 7 (NIPS 1994)}, 1995.
\bibitem[Wood et~al.(2023)]{wood2023unified} D.~Wood, T.~Mu, A.~Webb, H.~Reeve, M.~Luj\'an, and G.~Brown. A unified theory of diversity in ensemble learning. \emph{JMLR}, 24(359), 2023.
\bibitem[Jiang et~al.(2017)]{jiang2017generalized} Z.~Jiang, H.~Liu, B.~Fu, and Z.~Wu. Generalized ambiguity decompositions for classification with applications in active learning and unsupervised ensemble pruning. \emph{AAAI}, 2017.
\bibitem[Kuncheva and Whitaker(2003)]{kuncheva2003measures} L.~I.~Kuncheva and C.~J.~Whitaker. Measures of diversity in classifier ensembles and their relationship with the ensemble accuracy. \emph{Machine Learning}, 51(2):181-207, 2003.
\bibitem[Wang et~al.(2023)]{wang2023selfconsistency} X.~Wang, J.~Wei, D.~Schuurmans, Q.~Le, E.~Chi, S.~Narang, A.~Chowdhery, and D.~Zhou. Self-consistency improves chain of thought reasoning in language models. \emph{ICLR}, 2023.
\bibitem[Chen et~al.(2023)]{chen2023universal} X.~Chen et~al. Universal self-consistency for large language model generation. arXiv:2311.17311, 2023.
\bibitem[Wang et~al.(2025)]{wang2024moa} J.~Wang, J.~Wang, B.~Athiwaratkun, C.~Zhang, and J.~Zou. Mixture-of-Agents enhances large language model capabilities. \emph{ICLR}, 2025.
\bibitem[Li et~al.(2025)]{li2025rethinking} W.~Li et~al. Rethinking Mixture-of-Agents: is mixing different large language models beneficial? arXiv:2502.00674, 2025.
\bibitem[Naik et~al.(2023)]{naik2023diversity} R.~Naik, V.~Chandrasekaran, M.~Yuksekgonul, H.~Palangi, and B.~Nushi. Diversity of thought improves reasoning abilities of LLMs. arXiv:2310.07088, 2023.
\bibitem[Dipper(2024)]{dipper2024} Dipper: Diversity in prompts for producing large language model ensembles in reasoning tasks. arXiv:2412.15238, 2024.
\bibitem[Gao et~al.(2023)]{gao2023pal} L.~Gao, A.~Madaan, S.~Zhou, U.~Alon, P.~Liu, Y.~Yang, J.~Callan, and G.~Neubig. PAL: Program-aided language models. \emph{ICML}, 2023.
\bibitem[Chen et~al.(2022)]{chen2022program} W.~Chen, X.~Ma, X.~Wang, and W.~W.~Cohen. Program of thoughts prompting: disentangling computation from reasoning for numerical reasoning tasks. \emph{TMLR}, 2023. arXiv:2211.12588.
\bibitem[Liu et~al.(2023)]{liu2023plan} T.~Liu, Q.~Guo, Y.~Yang, X.~Hu, Y.~Zhang, X.~Qiu, and Z.~Zhang. Plan, verify and switch: integrated reasoning with diverse X-of-thoughts. \emph{EMNLP}, 2023.
\bibitem[Liang et~al.(2024)]{liang2024improving} Z.~Liang et~al. Improving LLM reasoning through scaling inference computation with collaborative verification. arXiv:2410.05318, 2024.
\bibitem[Audit(2026)]{audit2026diversity} Are diversity metrics measuring diversity? A capability-controlled audit of majority-vote gain in LLM ensembles. arXiv:2607.20768, 2026.
\bibitem[Balunovi\'c et~al.(2025)]{balunovic2025matharena} M.~Balunovi\'c, J.~Dekoninck, I.~Petrov, N.~Jovanovi\'c, and M.~Vechev. MathArena: evaluating LLMs on uncontaminated math competitions. arXiv:2505.23281, 2025.
\end{thebibliography}
\end{document}
